% Generated draft: requires mathematical and source review.
\documentclass[11pt]{article}
\usepackage[margin=25mm]{geometry}
\usepackage{fontspec}
\setmainfont{Latin Modern Roman}
\usepackage{amsmath,amssymb,mathtools,unicode-math}
\setmathfont{Latin Modern Math}
\usepackage{enumitem,needspace,xurl,hyperref}
\setlength{\parindent}{0pt}
\setlength{\parskip}{5pt}
\setlength{\emergencystretch}{3em}
\newcommand{\N}{\mathbb{N}}
\newcommand{\Z}{\mathbb{Z}}
\newcommand{\Q}{\mathbb{Q}}
\newcommand{\R}{\mathbb{R}}
\newcommand{\C}{\mathbb{C}}
\newcommand{\F}{\mathbb{F}}
\newcommand{\A}{\mathbb{A}}
\newcommand{\PP}{\mathbb{P}}
\newcommand{\E}{\mathbb{E}}
\newcommand{\eps}{\varepsilon}
\newcommand{\dd}{\,\mathrm{d}}
\begin{document}
\section*{UnsolvedMath ID: 30002341}
\textbf{Source problem code:} catalogue entry
\subsection{Definitions and mathematical statement}
Let \(F\) be a nonarchimedean local field and \(\mathcal O_F\) its ring of integers. The source considers a homogeneous spherical \(G\)-scheme \(X\) over \(\mathcal O_F\). Write \(G_F\) and \(X_F\) for their generic fibres. In the usual reductive setting, geometric homogeneity means that \(G_{\overline F}\) acts transitively on \(X_{\overline F}\), and sphericity means that a Borel subgroup of the connected reductive group \(G_{\overline F}\) has an open orbit on \(X_{\overline F}\). These generic-fibre conditions explain the terminology; the source does not specify the required integral-model conditions or the precise hypotheses on \(G\).

To make the requested algebra and module explicit, the following is a conditional formulation under the conventional interpretation using \(F\)-points and complex-valued functions. Set \(L=G_F(F)\) and \(Y=X_F(F)\), and let \(K\subseteq L\) be a compact open subgroup. Normalize a left Haar measure \(\dd g\) on \(L\) by \(\operatorname{vol}(K)=1\). Define
\[
\mathcal H(L,K)=\{h\in C_c^\infty(L,\C):h(k_1gk_2)=h(g)\text{ for all }g\in L,\ k_1,k_2\in K\},
\]
where \(C_c^\infty\) denotes locally constant, compactly supported functions. Its multiplication is convolution:
\[
(h_1*h_2)(g)=\int_L h_1(a)h_2(a^{-1}g)\,\dd a.
\]
For a left action of \(L\) on \(Y\), define
\[
M(Y,K)=\{f\in C_c^\infty(Y,\C):f(k^{-1}y)=f(y)\text{ for all }k\in K,\ y\in Y\}.
\]
The associated left Hecke action is
\[
(h\cdot f)(y)=\int_L h(g)f(g^{-1}y)\,\dd g.
\]
These formulas specify conventions for the objects to be described; they do not supply the requested structural description.

\textbf{Problem.} For every homogeneous spherical \(G\)-scheme \(X\) over \(\mathcal O_F\) within the source's intended class, describe the spherical Hecke algebra \(\mathcal H(G,K)\) and the module structure of \(C_c^\infty(X)^K\) over that algebra. Under the conditional conventions above, this asks for a structural description of \(\mathcal H(L,K)\) together with its action on \(M(Y,K)\).

\textbf{Unresolved qualifications.} The record does not identify \(K\), the coefficient field, the point space denoted by \(X\) in the function notation, or the exact reductive and integral hypotheses. In particular, it does not state whether \(K=G(\mathcal O_F)\), whether this subgroup is hyperspecial, or whether splitness or unramifiedness is required. The conditional formulation therefore requires confirmation before it can be treated as an exact formulation of the source problem. The supplied literature assessment mentions a uniform basis, whereas the problem statement itself requests a module description without prescribing a basis.
\subsection{Short English statement}
For a homogeneous spherical scheme over the ring of integers of a nonarchimedean local field, describe the spherical Hecke algebra and its action on the \(K\)-invariant, locally constant, compactly supported functions. The intended subgroup \(K\), function space, coefficients, and integral hypotheses need confirmation. The supplied record labels the problem partially solved; this is a dated curation claim, not an independently verified conclusion.
\subsection{Sources}
\begin{itemize}
\item Supplied local UnsolvedMath registry record, permanent ID \(30002341\), title \emph{Spherical Hecke Modules on Homogeneous Schemes}, collection \emph{Oberwolfach Reports Open Problems}, exact source problem code \texttt{catalogue entry}. No individual proposer is supplied; the proposed year is recorded as 2013.
\item Supplied source citation: \emph{Mini-Workshop: Spherical Varieties and Automorphic Representations (2014)}, Oberwolfach Reports, pp.\ 1445--1494, DOI \texttt{10.4171/owr/2013/24}. Supplied source URL and literature reference, labeled \emph{Oberwolfach Report 2013/24}: \url{https://doi.org/10.4171/owr/2013/24}. The citation's 2014 date and the record's 2013 identifiers are retained without reconciliation.
\item Supplied literature reference: Sakellaridis--Venkatesh, \emph{Periods and harmonic analysis on spherical varieties}, \url{https://arxiv.org/abs/1203.0039}.
\item Literature assessment dated exactly \texttt{2026-08-21}: the supplied curation labels the problem \emph{partially solved} and states that generalized Satake transforms and substantial spherical-variety cases describe the Hecke-module structure, while a uniform basis and module description for arbitrary spherical \(X\) remain incomplete. This is supplied dated curation, not independently verified literature assessment. The background expressly requests independent verification of time-sensitive status and cited claims.
\item The registry labels the statement \texttt{corrected\_verified} and reports that its mathematical content was checked against the official Oberwolfach report, with notation, context, and question form normalized without changing the source-level claim. This verification claim is supplied metadata; no independent source consultation is asserted here.
\end{itemize}
\end{document}
